\documentclass{article}
\usepackage{lmodern}
\usepackage{amsmath}
\title{From Mathematics to Generic Programming}
\author{Companion repository: corrected solution}
\date{November 2026}
\begin{document}
\maketitle
\section*{Exercise 4.4: Odd and even squares}
\textit{Proposition.} For any odd square $x$, there is an even square $y$
such that $x+y$ is a square.

\textit{Proof.} Start with $9+16=25$, or $3^2+4^2=5^2$.
We need an even \emph{square}, not merely an even number.
Write $x=(2k+1)^2=4k(k+1)+1$. Since one of the consecutive integers
$k$ and $k+1$ is even, $x\equiv1\pmod8$.
Consequently $m=(x-1)/2$ is even (indeed, a multiple of four).
Set
\[
 y=m^2=\left(\frac{x-1}{2}\right)^2.
\]
Then $y$ is an even square and, since $x=2m+1$,
\[
 x+y=2m+1+m^2=(m+1)^2.
\]

\paragraph{Boundary case.}
For $x=1$, this construction gives $m=0$ and $y=0=0^2$, so
$1+0=1^2$. The exercise does not explicitly say ``positive squares.''
Under the convention that squares include $0^2$, the proposition includes
this case. If squares are required to be positive, the proposition must
instead assume $x>1$: two positive integer squares cannot differ by one,
so no positive square $y$ works when $x=1$.
\end{document}
